% Shared preamble for the Carina theory manual.
%
% Every chapter writes tensors and matrices through the macros defined here,
% so the notation is set in one place and cannot drift between chapters.

\usepackage[T1]{fontenc}
\usepackage{lmodern}
\usepackage{amsmath, amssymb, amsthm}
\usepackage{boldtensors}
\usepackage{booktabs}
\usepackage{microtype}
% `algorithmic` rather than `algpseudocode`: the latter lives in algorithmicx,
% which is not in a base TeX Live installation.  This manual builds against
% texlive-latex-recommended with no extra package installs.
\usepackage{algorithm}
\usepackage{algorithmic}
\usepackage[margin=1.15in]{geometry}
\usepackage[hidelinks]{hyperref}

% File paths set in typewriter cannot hyphenate, so a paragraph containing one
% can overflow the measure.  Let TeX stretch the interword space as a last
% resort rather than run into the margin.
\setlength{\emergencystretch}{2.5em}

% --------------------------------------------------------------------------
% Notation
% --------------------------------------------------------------------------
% boldtensors makes \bf-prefixed Greek and upright symbols bold, which is what
% carries the tensor/matrix distinction throughout.  Second-order tensors and
% matrices are bold upright capitals, vectors bold lowercase, fourth-order
% tensors blackboard bold.

% boldtensors makes `~` and `"` active in math mode: `~X` sets X in the bold
% tensor font, `"X` in blackboard bold.  That is the package's own interface;
% wrapping it in named macros keeps the active characters out of the chapter
% sources, where they would be easy to mistype and hard to grep for.

\newcommand{\tensor}[1]{~#1}                  % 2nd-order tensor / matrix
\newcommand{\vect}[1]{~#1}                    % vector
\newcommand{\fourth}[1]{"#1}                  % 4th-order tensor

% Continuum quantities
\newcommand{\Fdef}{\tensor{F}}                % deformation gradient
\newcommand{\Piola}{\tensor{P}}               % first Piola--Kirchhoff stress
\newcommand{\Cauchy}{\tensor{\sigma}}         % Cauchy stress
\newcommand{\Kirch}{\tensor{\tau}}            % Kirchhoff stress
\newcommand{\strain}{\tensor{\varepsilon}}    % small strain
\newcommand{\Cright}{\tensor{C}}              % right Cauchy--Green
\newcommand{\Ident}{\tensor{I}}               % identity tensor
\newcommand{\disp}{\vect{u}}                  % displacement field
\newcommand{\Grad}{\nabla}

% Discrete quantities
\newcommand{\Kmat}{\tensor{K}}                % tangent stiffness
\newcommand{\Mmat}{\tensor{M}}                % mass matrix
\newcommand{\Keff}{\tensor{K}_{\mathrm{eff}}} % Newmark effective stiffness
\newcommand{\Rvec}{\vect{R}}                  % residual
\newcommand{\Uvec}{\vect{u}}                  % discrete displacement
\newcommand{\Vvec}{\vect{v}}                  % discrete velocity
\newcommand{\Avec}{\vect{a}}                  % discrete acceleration
\newcommand{\Fint}{\vect{f}^{\mathrm{int}}}
\newcommand{\Fext}{\vect{f}^{\mathrm{ext}}}
\newcommand{\Pmat}{\tensor{P}}                % prolongation
\newcommand{\Rrest}{\tensor{R}}               % restriction
\newcommand{\Dmat}{\tensor{D}}                % diagonal of an operator
\newcommand{\Bnull}{\tensor{B}}               % near-nullspace
\newcommand{\Prec}{\tensor{M}}                % preconditioner

% Operators
\DeclareMathOperator{\dev}{dev}
\DeclareMathOperator{\tr}{tr}
\DeclareMathOperator{\diag}{diag}
\DeclareMathOperator*{\argmin}{arg\,min}
\newcommand{\norm}[1]{\left\lVert #1 \right\rVert}
\newcommand{\abs}[1]{\left\lvert #1 \right\rvert}
\newcommand{\inner}[2]{\left\langle #1, #2 \right\rangle}
\newcommand{\dbldot}{\mathbin{:}}

% Cross-references into the implementation.  Every chapter names the files that
% implement it, so the two can be checked against each other.
%
% Set in the text flow rather than as a \marginpar: the margins here are too
% narrow for a path like `src/nonlinear_solvers.jl`, and a marginpar that does
% not fit is silently clipped at the page edge rather than reported.
\newcommand{\impl}[1]{%
  \vspace{-0.6em}%
  {\raggedleft\small\textit{Implementation:}~\texttt{#1}\par}%
  \vspace{0.4em}%
}
\newcommand{\code}[1]{\texttt{#1}}

\theoremstyle{definition}
\newtheorem{remark}{Remark}[chapter]
\theoremstyle{plain}
\newtheorem{proposition}{Proposition}[chapter]
